
\documentclass{article}
\usepackage{colortbl}
\usepackage{makecell}
\usepackage{multirow}
\usepackage{supertabular}

\begin{document}

\newcounter{utterance}

\centering \large Interaction Transcript for game `referencegame', experiment `letter\_grids', episode 16 with reflexion{-}Qwen3.5{-}4B.
\vspace{24pt}

{ \footnotesize  \setcounter{utterance}{1}
\setlength{\tabcolsep}{0pt}
\begin{supertabular}{c@{$\;$}|p{.15\linewidth}@{}p{.15\linewidth}p{.15\linewidth}p{.15\linewidth}p{.15\linewidth}p{.15\linewidth}}
    \# & $\;$A & \multicolumn{4}{c}{Game Master} & $\;\:$B\\
    \hline

    \theutterance \stepcounter{utterance}  
    & & \multicolumn{4}{p{0.6\linewidth}}{
        \cellcolor[rgb]{0.9,0.9,0.9}{
            \makecell[{{p{\linewidth}}}]{
                \texttt{\tiny{[P1$\langle$GM]}}
                \texttt{You have attempted this same task before and did not succeed. These are the plans you wrote after those attempts.} \\
\texttt{Trial \#1: New plan: My previous description was too verbose and listed every single filled cell individually, which likely made the expression ambiguous or failed to capture the distinctive pattern that separates the target from the distractors. The target grid is unique because it has a "U" shape formed by the top and bottom rows, with only three specific cells filled in the middle rows (one in row 2, one in row 3, one in row 4). The distractors have different internal patterns (alternating or blocky). Next time, I will construct a concise expression that focuses on the global structure (e.g., "full top and bottom rows") and the specific, sparse arrangement of the middle cells without listing coordinates for every single one, ensuring the description uniquely identifies the target's specific "U" shape with three internal points.} \\
\\ 
\\ 
\texttt{You are given three grids, where each of them is 5 by 5 in size.} \\
\texttt{Grids have empty cells marked with "▢" and filled cells marked with "X".} \\
\texttt{Your task is to generate a referring expression that best describes the target grid while distinguishing it from the two other distractor grids.} \\
\texttt{The first grid is the target grid, and the following two grids are the distractors.} \\
\\ 
\texttt{Target grid:} \\
\\ 
\texttt{X X X X X} \\
\texttt{□ □ □ X X} \\
\texttt{□ □ X X □} \\
\texttt{□ X X □ □} \\
\texttt{X X X X X} \\
\\ 
\texttt{Distractor grid 1:} \\
\\ 
\texttt{X X X X X} \\
\texttt{X □ X □ X} \\
\texttt{X X X X X} \\
\texttt{X □ X □ X} \\
\texttt{X X X X X} \\
\\ 
\texttt{Distractor grid 2:} \\
\\ 
\texttt{X X X X X} \\
\texttt{X □ □ □ □} \\
\texttt{X □ □ □ □} \\
\texttt{X □ □ □ □} \\
\texttt{X X X X X} \\
\\ 
\texttt{Instruction: Describe the target grid.} \\
\texttt{Generate the referring expression starting with the tag "Expression: " for the given target grid. Omit any other text.} \\
            }
        }
    }
    & \\ \\

    \theutterance \stepcounter{utterance}  
    & \multicolumn{4}{p{0.6\linewidth}}{
        \cellcolor[rgb]{0.9,0.9,0.9}{
            \makecell[{{p{\linewidth}}}]{
                \texttt{\tiny{[P1$\rangle$GM]}}
                \texttt{Expression: A grid with full top and bottom rows, containing exactly three filled cells in the middle rows: one in the fourth column of the second row, two adjacent filled cells in the third and fourth columns of the third row, and two adjacent filled cells in the second and third columns of the fourth row.} \\
            }
        }
    }
    & & \\ \\

    \theutterance \stepcounter{utterance}  
    & & & \multicolumn{2}{p{0.3\linewidth}}{
        \cellcolor[rgb]{0.9,0.9,0.9}{
            \makecell[{{p{\linewidth}}}]{
                \texttt{\tiny{[GM$|$GM]}}
                \texttt{Expression: A grid with full top and bottom rows, containing exactly three filled cells in the middle rows: one in the fourth column of the second row, two adjacent filled cells in the third and fourth columns of the third row, and two adjacent filled cells in the second and third columns of the fourth row.} \\
            }
        }
    }
    & & \\ \\

    \theutterance \stepcounter{utterance}  
    & & \multicolumn{4}{p{0.6\linewidth}}{
        \cellcolor[rgb]{0.9,0.9,0.9}{
            \makecell[{{p{\linewidth}}}]{
                \texttt{\tiny{[GM$\rangle$P2]}}
                \texttt{You have attempted this same task before and did not succeed. These are the plans you wrote after those attempts.} \\
\texttt{Trial \#1: New plan: I incorrectly identified the target grid by relying on a partial match of the description rather than verifying every specific condition. The expression requires a filled cell in the second row at the fourth column, but the First grid has an empty cell there. Next time, I will systematically check each grid against all constraints in the expression: full top/bottom rows, the specific single filled cells in rows 2, 3, and 4 at columns 4, 3, and 2 respectively, ensuring no single condition is missed before selecting an answer.} \\
\\ 
\\ 
\texttt{You are given three grids, where each of them is 5 by 5 in size.} \\
\texttt{Grids have empty cells marked with "▢" and filled cells marked with "X".} \\
\texttt{You are also given a referring expression that describes one of the given grids.} \\
\texttt{Your task is to select the grid that matches the given referring expression.} \\
\texttt{Generate only the number (in text) of the grid that matches the given expression by selecting first, second, or third.} \\
\\ 
\texttt{First:} \\
\\ 
\texttt{X X X X X} \\
\texttt{X □ X □ X} \\
\texttt{X X X X X} \\
\texttt{X □ X □ X} \\
\texttt{X X X X X} \\
\\ 
\texttt{Second:} \\
\\ 
\texttt{X X X X X} \\
\texttt{□ □ □ X X} \\
\texttt{□ □ X X □} \\
\texttt{□ X X □ □} \\
\texttt{X X X X X} \\
\\ 
\texttt{Third:} \\
\\ 
\texttt{X X X X X} \\
\texttt{X □ □ □ □} \\
\texttt{X □ □ □ □} \\
\texttt{X □ □ □ □} \\
\texttt{X X X X X} \\
\\ 
\texttt{Expression: A grid with full top and bottom rows, containing exactly three filled cells in the middle rows: one in the fourth column of the second row, two adjacent filled cells in the third and fourth columns of the third row, and two adjacent filled cells in the second and third columns of the fourth row.} \\
\texttt{Question: Which grid does the expression refer to?} \\
\texttt{Start with the tag "Answer: ", followed by your selection. Omit any other text.} \\
            }
        }
    }
    & \\ \\

    \theutterance \stepcounter{utterance}  
    & & & \multicolumn{4}{p{0.6\linewidth}}{
        \cellcolor[rgb]{0.9,0.9,0.9}{
            \makecell[{{p{\linewidth}}}]{
                \texttt{\tiny{[GM$\langle$P2]}}
                \texttt{Answer: Third} \\
            }
        }
    }
     \\ \\

    \theutterance \stepcounter{utterance}  
    & & & \multicolumn{2}{p{0.3\linewidth}}{
        \cellcolor[rgb]{0.9,0.9,0.9}{
            \makecell[{{p{\linewidth}}}]{
                \texttt{\tiny{[GM$|$GM]}}
                \texttt{third} \\
            }
        }
    }
    & & \\ \\

    \theutterance \stepcounter{utterance}  
    & \multicolumn{4}{p{0.6\linewidth}}{
        \cellcolor[rgb]{0.95,0.95,0.95}{
            \makecell[{{p{\linewidth}}}]{
                \texttt{\tiny{[P2$\rightarrow$P1]}}
                \texttt{Post{-}episode Reflexion (not sent to GM):} \\
\texttt{New plan: My previous description was too detailed and listed specific coordinates for every filled cell, which likely caused the system to reject it as non{-}concise or failed to capture the high{-}level distinguishing feature. The target grid is uniquely defined by having exactly one filled cell in the second row (column 4) and exactly one filled cell in the fourth row (column 2), whereas the distractors have two or more filled cells in those middle rows. Next time, I will construct a concise expression that focuses on the global structure (full top and bottom rows) and the specific count and position of the single filled cells in the second and fourth rows, omitting the third row's details entirely to ensure uniqueness.} \\
            }
        }
    }
    & & \\ \\

    \theutterance \stepcounter{utterance}  
    & & & \multicolumn{4}{p{0.6\linewidth}}{
        \cellcolor[rgb]{0.95,0.95,0.95}{
            \makecell[{{p{\linewidth}}}]{
                \texttt{\tiny{[P1$\rightarrow$P2]}}
                \texttt{Post{-}episode Reflexion (not sent to GM):} \\
\texttt{New plan: I incorrectly identified the target grid by relying on a partial match of the description rather than verifying every specific condition. The expression requires a filled cell in the second row at the fourth column, but the Third grid has an empty cell there. Next time, I will systematically check each grid against all constraints in the expression: full top/bottom rows, the specific single filled cell in row 2 column 4, the two adjacent filled cells in row 3 columns 3{-}4, and the two adjacent filled cells in row 4 columns 2{-}3, ensuring no single condition is missed before selecting an answer.} \\
            }
        }
    }
     \\ \\

\end{supertabular}
}

\end{document}
